\refstepcounter{subsection}\label{X.2.6}
\begin{enonce*}{Corollaire 2.6\ndemark}
Si on~a $\Lef(X, Y)$ et si $Y$ est connexe, tout voisinage ouvert $U$ de $Y$ est connexe et l'homomorphisme naturel $\pi_1(Y)\to\pi_1(U)$ est surjectif. Si de plus on~a $\Leff(X, Y)$, l'homomorphisme naturel
\[
\pi_1(Y) \to \varprojlim_{U}\, \pi_1(U)
\]
est un isomorphisme. (N.B. On suppose choisi un \og point-base\fg dans $Y$, qu'on prend aussi pour point-base dans $X$, pour la d\'efinition des groupes fondamentaux.)
\end{enonce*}
